\chapter[25]{Estimates of the Heat Equation for Evolving Metrics}
\label{ch:chow-iii-estimates-heat-equation-evolving-metrics}
\dcref{ch25}
\source{chow-et-al-ricci-flow-part-iii:page-0326}

Let \(g(\tau)\), \(\tau\in[0,T]\), be a smooth family of complete metrics
with
\[
 \partial_\tau g_{ij}=2\mathcal R_{ij},\qquad
 \mathcal R=g^{ij}\mathcal R_{ij}.
\]
We consider positive subsolutions and solutions of
\[
 \partial_\tau u\leq \Delta_{g(\tau)}u-Qu.
\]

\section{The parabolic mean value inequality}

\begin{definition}[Parabolic mean value inequality]
\label{def:chow-iii-parabolic-mean-value-inequality}
\source{chow-et-al-ricci-flow-part-iii:page-0326,chow-et-al-ricci-flow-part-iii:page-0327}
\dcref{ch25:25.1}
\group{Parabolic Mean Value Inequalities}

A complete Riemannian manifold \((\mathcal M^n,g)\) satisfies the
\emph{parabolic mean value inequality} up to time \(T\) with constant \(C\)
if every nonnegative subsolution \(f\) of \(\partial_t f\leq\Delta f\) on
\(P=B_g(p,\sqrt t)\times[0,t]\) satisfies
\[
 f(p,t)\leq \frac{C}{\Vol_{g+dt^2}(P)}
 \int_0^t\int_{B_g(p,\sqrt t)}f\,d\mu_g\,ds.
\]
\end{definition}

Fix a complete background metric \(\widetilde g\), a compact smooth domain
\(\Omega\), and assume
\[
 C_0^{-1}\widetilde g\leq g(0)\leq C_0\widetilde g,
 \qquad
 \Lambda=\sup_{\Omega\times[0,T]}|\mathcal R_{ij}|_{g(\tau)}<\infty.
\]
Then
\[
 \widetilde C_0^{-1}\widetilde g\leq g(\tau)
 \leq\widetilde C_0\widetilde g,
 \qquad \widetilde C_0=C_0e^{2\Lambda T}.
\]
Write
\[
 \mathcal P_{\widetilde g}(x,\tau,r,-r^2)
 =B_{\widetilde g}(x,r)\times[\tau-r^2,\tau].
\]

\begin{theorem}[Parabolic mean value inequality]
\label{thm:chow-iii-evolving-parabolic-mean-value}
\source{chow-et-al-ricci-flow-part-iii:page-0327,chow-et-al-ricci-flow-part-iii:page-0328,chow-et-al-ricci-flow-part-iii:page-0337}
\uses{lem:chow-iii-reverse-poincare, prop:chow-iii-local-sobolev-ball, lem:chow-iii-moser-cylinder-iteration}
\dcref{ch25:25.2}
\group{Parabolic Mean Value Inequalities}


Suppose \(\Ric(\widetilde g)\geq-K\) on \(\Omega\), and let \(u>0\)
satisfy
\[
 \partial_\tau u\leq\Delta_{g(\tau)}u-Qu.
\]
If
\[
 \overline{\mathcal P}_{\widetilde g}
 (x_0,\tau_0,4r_0,-(2r_0)^2)\subset\Omega\times[0,T],
\]
then
\[
 \sup_{\mathcal P_{\widetilde g}(x_0,\tau_0,r_0,-r_0^2)}u
 \leq
 \frac{C_1e^{C_2\tau_0+C_3\sqrt K r_0}}
 {r_0^2\Vol_{\widetilde g}(B_{\widetilde g}(x_0,r_0))}
 \int_{\mathcal P_{\widetilde g}(x_0,\tau_0,2r_0,-(2r_0)^2)}
 u\,d\mu_{\widetilde g}\,d\tau.
\]
Here \(C_1\) depends only on \(C_0,n,T,\Lambda\), \(C_2\) only on
\(\sup|Q|\) and \(\sup|\mathcal R|\), and \(C_3\) only on \(n\).
\end{theorem}
\begin{proof}
\sketch Choose \(A\geq0\) so that \(Q+A-\mathcal R/(2p)\geq0\) for all
\(p\geq1\), and put \(v=e^{-A\tau}u\).  Lemma
\ref{lem:chow-iii-reverse-poincare}, followed by Proposition
\ref{prop:chow-iii-local-sobolev-ball}, gives the one-step estimate
\[
 \|v\|_{L^{2p(n+2)/n}(P_{r'})}
 \leq M(r,r',p)^{1/(2p)}\|v\|_{L^{2p}(P_r)}.
\]
Apply Lemma~\ref{lem:chow-iii-moser-cylinder-iteration} on radii decreasing
from \(2r_0\) to \(r_0\).  This first bounds \(L^\infty(P_{r_0})\) by
\(L^2(P_{2r_0})\).  On radii increasing from \(r_0\) to \(2r_0\), use
\(\|v\|_2\leq\|v\|_1^{1/2}\|v\|_\infty^{1/2}\) and iterate once more.
The geometric exponents sum to finite constants, yielding
\[
 \|v\|_{L^\infty(P_{r_0})}
 \leq \frac{C e^{C\sqrt K r_0}}
 {r_0^2\Vol_{\widetilde g}(B_{\widetilde g}(x_0,r_0))}
 \|v\|_{L^1(P_{2r_0})}.
\]
Finally \(v\leq u\leq e^{A\tau_0}v\) on the cylinder.  The stated
dependences follow from the choice of \(A\), the metric comparison, and the
local Sobolev constant.
\end{proof}

\begin{lemma}[Reverse Poincare-type inequalities]
\label{lem:chow-iii-reverse-poincare}
\source{chow-et-al-ricci-flow-part-iii:page-0328,chow-et-al-ricci-flow-part-iii:page-0329,chow-et-al-ricci-flow-part-iii:page-0330}
\dcref{ch25:25.3}
\group{Moser Iteration}

Let \(v=e^{-A\tau}u\), where
\(Q+A-\mathcal R/(2p)\geq0\) for every \(p\geq1\).  Suppose
\(0\leq\psi\leq1\) has compact spatial support in \(D\Subset\Omega\),
\(\psi(\cdot,\tau_1)=0\), and
\[
 \psi\partial_\tau\psi+|\nabla\psi|_{g(\tau)}^2\leq L.
\]
Then
\[
 \int_{\tau_1}^{\tau_2}\int_D|\nabla(\psi v^p)|^2\,d\mu\,d\tau
 \leq L\int_{\tau_1}^{\tau_2}\int_Dv^{2p}\,d\mu\,d\tau
\]
and
\[
 \left(\int_D\psi^2v^{2p}\,d\mu\right)(\tau_2)
 \leq2L\int_{\tau_1}^{\tau_2}\int_Dv^{2p}\,d\mu\,d\tau.
\]
\end{lemma}
\begin{proof}
\sketch The subsolution inequality gives
\[
 \partial_\tau(v^p)-\Delta(v^p)+p(Q+A)v^p\leq0.
\]
Multiply it by \(\psi^2v^p\), integrate over space-time, and use
\(\partial_\tau d\mu=\mathcal R\,d\mu\).  Spatial integration by parts gives
\[
 -\int_D\psi^2v^p\Delta(v^p)
 =\int_D|\nabla(\psi v^p)|^2-\int_Dv^{2p}|\nabla\psi|^2.
\]
The initial boundary term vanishes, while the terminal term is
\(\frac12\int_D\psi^2v^{2p}\).  The choice of \(A\) makes the remaining
zeroth-order term nonnegative.  Bounding the two cutoff terms by \(Lv^{2p}\)
and then discarding either nonnegative term proves the two inequalities.
\end{proof}

\begin{proposition}[Local \(L^2\)-Sobolev inequality in a ball]
\label{prop:chow-iii-local-sobolev-ball}
\source{chow-et-al-ricci-flow-part-iii:page-0330}
\dcref{ch25:25.4}
\group{Moser Iteration}

There is \(C_S=C_S(n)\) such that, if \(n\geq3\) and
\(\Ric_{\widetilde g}\geq-K\) on \(B_{\widetilde g}(x_0,2r)\), then every
smooth compactly supported \(f\) on \(B_{\widetilde g}(x_0,r)\) satisfies
\[
 \left(\int |f|^{2n/(n-2)}\,d\mu_{\widetilde g}\right)^{(n-2)/n}
 \leq
 \frac{e^{C_S(1+\sqrt K r)}}
 {\Vol_{\widetilde g}(B_{\widetilde g}(x_0,r))^{2/n}}
 \int\bigl(r^2|\nabla f|^2+f^2\bigr)\,d\mu_{\widetilde g}.
\]
\end{proposition}
\begin{proof}
\sketch Rescale to \(r=1\).  Bishop--Gromov comparison under
\(\Ric\geq-K\) gives a relative isoperimetric inequality on the unit ball,
with loss at most \(e^{C(n)(1+\sqrt K)}\) and normalization by the ball's
volume.  Applying that inequality to the level sets of \(|f|\), integrating
with the coarea formula, and then applying Holder yields the normalized
\(W^{1,1}\) Sobolev inequality.  Apply it to
\(|f|^{2(n-1)/(n-2)}\), use Cauchy--Schwarz on the gradient term, and absorb
the resulting lower-order term.  Scaling back restores \(r^2\) and the
factor \(\Vol(B(x_0,r))^{-2/n}\).
\end{proof}

\begin{remark}[The two-dimensional Sobolev exponent]
\label{rem:chow-iii-two-dimensional-local-sobolev}
\source{chow-et-al-ricci-flow-part-iii:page-0331}
\dcref{ch25:25.5}
\group{Moser Iteration}

When \(n=2\), for every \(q>2\) the preceding estimate holds with target
exponent \(2q/(q-2)\), volume exponent \(-2/q\), and a constant depending
only on \(q\).
\end{remark}

\begin{lemma}[Moser iteration on nested cylinders]
\label{lem:chow-iii-moser-cylinder-iteration}
\source{chow-et-al-ricci-flow-part-iii:page-0331,chow-et-al-ricci-flow-part-iii:page-0332,chow-et-al-ricci-flow-part-iii:page-0333,chow-et-al-ricci-flow-part-iii:page-0334,chow-et-al-ricci-flow-part-iii:page-0335}
\uses{lem:chow-iii-reverse-poincare, prop:chow-iii-local-sobolev-ball}
\dcref{ch25:25.6}
\group{Moser Iteration}


For \(r\in[r_0,2r_0)\) and \(1<\gamma\leq2\), with
\(\gamma r\leq2r_0\), one has
\[
 \|v\|_{L^\infty(P_r)}\leq \widehat C(\gamma)
 \|v\|_{L^2(P_{\gamma r})},
\]
where
\[
 \widehat C(\gamma)=
 \frac{e^{C_S(1+\gamma\sqrt K r)n/4}}
 {r\Vol_{\widetilde g}(B_{\widetilde g}(x_0,r))^{1/2}}
 \widetilde C_0^{(n+3)n/4}\widehat C_n
 (\gamma-1)^{-n/4}
\]
and \(\widehat C_n\) depends only on \(n\).
\end{lemma}
\begin{proof}
\sketch Put \(\beta=(n+2)/n\), \(p_i=\beta^{i-1}\), and
\[
 r_i=r\left(1+\frac{\gamma-1}{2^{i-1}}\right).
\]
Choose space-time cutoffs equal to one on \(P_{r_{i+1}}\), supported in
\(P_{r_i}\), with
\(\psi_i\partial_\tau\psi_i+|\nabla\psi_i|^2
 \leq C4^i((\gamma-1)r)^{-2}\).
Lemma~\ref{lem:chow-iii-reverse-poincare}, Proposition
\ref{prop:chow-iii-local-sobolev-ball}, and Holder give
\[
 \|v\|_{L^{2p_{i+1}}(P_{r_{i+1}})}
 \leq M_i^{1/(2p_{i+1})}
 \|v\|_{L^{2p_i}(P_{r_i})},
\]
where \(M_i\) is bounded by a fixed geometric factor times a power of
\(4^i((\gamma-1)r)^{-2}\).  Iteration is legitimate because
\(\sum\beta^{-i}<\infty\) and \(\sum i\beta^{-i}<\infty\).  Letting
\(i\to\infty\), and using volume comparison to replace all ball volumes by
that of \(B(x_0,r)\), gives the displayed constant.
\end{proof}

\begin{example}[Rescaled nested-cylinder iteration]
\label{ex:chow-iii-rescaled-moser-cylinder-iteration}
\source{chow-et-al-ricci-flow-part-iii:page-0335}
\uses{lem:chow-iii-moser-cylinder-iteration}
\dcref{ch25:25.7}
\group{Moser Iteration}


For the radii
\(r_i=r(1+(\gamma-1)2^{1-i})\).
\end{example}
\begin{proof}
\sketch For these radii,
\(r_i-r_{i+1}=r(\gamma-1)2^{-i}\).  Spatial and temporal cutoffs therefore
have squared-gradient and time-derivative bounds
\(C4^i((\gamma-1)r)^{-2}\).  Substitution in the one-step Sobolev estimate
and multiplication over \(i\) gives powers
\(\sum\beta^{-i}=n/2\) and \(\sum i\beta^{-i}<\infty\).  The radius gap
contributes exactly \((\gamma-1)^{-n/4}\), proving the lemma.
\end{proof}

\section{Li--Yau differential Harnack estimates}

\begin{theorem}[Harnack estimate of Li--Yau type for evolving metrics]
\label{thm:chow-iii-li-yau-evolving-metrics}
\source{chow-et-al-ricci-flow-part-iii:page-0339,chow-et-al-ricci-flow-part-iii:page-0343,chow-et-al-ricci-flow-part-iii:page-0344,chow-et-al-ricci-flow-part-iii:page-0345,chow-et-al-ricci-flow-part-iii:page-0346,chow-et-al-ricci-flow-part-iii:page-0347,chow-et-al-ricci-flow-part-iii:page-0348,chow-et-al-ricci-flow-part-iii:page-0349,chow-et-al-ricci-flow-part-iii:page-0350,chow-et-al-ricci-flow-part-iii:page-0351}
\uses{cor:chow-iii-harnack-quantity-evolution-inequality, lem:bounded-curvature-distance-like-function}
\dcref{ch25:25.8}
\group{Li--Yau Estimates}


Assume on \(\mathcal M\times[0,T]\) that
\[
 \Ric\geq-K,\qquad |\mathcal R_{ij}|\leq\Lambda,
\]
\[
 \left|-2g^{ij}\nabla_i\mathcal R_{jk}+\nabla_k\mathcal R\right|
 +\frac43|\nabla Q|\leq C',\qquad \Delta Q\leq C',
\]
and \(|\sec(g(0))|\leq K\).  If \(R\geq1\) and \(u>0\) solves
\(\partial_\tau u=\Delta u-Qu\) on
\(B_{g(0)}(p,2R)\times[0,T]\), then for \(0<\varepsilon<2/3\),
\[
 \partial_\tau\log u-(1-\varepsilon)|\nabla\log u|^2+Q
 \geq-\frac{n}{2-3\varepsilon}
 \left(\frac1\tau+\frac{C_5}{R}+\frac{C_6}{R^2}+C_7\right)
\]
on \(B_{g(0)}(p,R-C_{n,K})\times(0,T]\).  The constants have the
dependencies stated in the source: \(C_5\) also depends on
\(\sup|\nabla\mathcal R|\), \(C_6\) on \(\varepsilon\), and \(C_7\) on
\(\varepsilon,C'\).  In particular, letting \(R\to\infty\) removes the
\(R^{-1}\) and \(R^{-2}\) terms.
\end{theorem}
\begin{proof}
\sketch Set \(L=\log u\) and
\(P=\Delta L+\varepsilon|\nabla L|^2\).  By Corollary
\ref{cor:chow-iii-harnack-quantity-evolution-inequality},
\[
 \partial_\tau P\geq\Delta P+2\langle\nabla L,\nabla P\rangle
 +\frac{2-3\varepsilon}{n}(\Delta L)^2-C_1|\nabla L|^2-C_2.
\]
Choose a smooth distance-like function \(f\) for \(g(0)\) with bounded first
and second derivatives, supplied by Lemma
\ref{lem:bounded-curvature-distance-like-function}, and put
\(\phi=\psi(f/R)\), where \(\psi=1\) on \([0,1]\), vanishes on
\([2,\infty)\), and satisfies
\(|\psi'|^2/\psi+|\psi''|\leq C\).  Metric comparison and the connection
variation formula give
\[
 |\nabla f|\leq C,qquad |\Delta_{g(\tau)}f|\leq C,
\]
with the required dependence on \(\sup|\nabla\mathcal R|\).  Consequently
\[
 (\partial_\tau-\Delta)\phi+
 \left(2+\frac{n}{2(2-3\varepsilon)\varepsilon}\right)
 \frac{|\nabla\phi|^2}{\phi}
 \leq C_5R^{-1}+C_6R^{-2}.
\]

At a negative minimum of \(\tau\phi P\), the first derivative vanishes and
the parabolic second-derivative inequality has the favorable sign.  Use
\(\Delta L=P-\varepsilon|\nabla L|^2\), and estimate the mixed cutoff term by
Peter--Paul with parameter
\(2(2-3\varepsilon)\varepsilon/n\).  After multiplying by \(\phi\), one gets
\[
 0\geq \frac{2-3\varepsilon}{n}(\phi P)^2
 +\left(C_3+\frac\phi\tau\right)\phi P
 -\phi^2\left(C_2+\frac{nC_1^2}
 {(2-3\varepsilon)\varepsilon^2}\right).
\]
Solving this quadratic inequality bounds \(\tau\phi P\) from below.
The cutoff equals one on \(B_{g(0)}(p,R-C_{n,K})\); substitution of the
constant bounds yields the asserted estimate.  Exhausting \(\mathcal M\)
by taking \(R\to\infty\) proves the global form.
\end{proof}

\begin{theorem}[Harnack estimate of Li--Yau type for Ricci flow]
\label{thm:chow-iii-li-yau-ricci-flow}
\source{chow-et-al-ricci-flow-part-iii:page-0339,chow-et-al-ricci-flow-part-iii:page-0340,chow-et-al-ricci-flow-part-iii:page-0351,chow-et-al-ricci-flow-part-iii:page-0352}
\uses{cor:chow-iii-harnack-quantity-evolution-inequality, thm:chow-iii-changing-distances}
\dcref{ch25:25.9}
\group{Li--Yau Estimates}


Let \(g(\tau)\) be a complete Ricci flow
\(\partial_\tau g=-2\Ric\).  Suppose that, on
\(B_{g(\tau)}(p,2R)\) for each \(\tau\),
\[
 |\Ric|\leq K,\qquad \frac43|\nabla Q|\leq C',
 \qquad\Delta Q\leq C',
\]
where \(R\geq\max\{1,K^{-1/2}\}\).  If \(u>0\) solves
\(\partial_\tau u=\Delta u-Qu\) there, then on
\(B_{g(\tau)}(p,R)\),
\[
 \partial_\tau\log u-(1-\varepsilon)|\nabla\log u|^2+Q
 \geq-\frac{n}{2-3\varepsilon}
 \left(\frac1\tau+\frac{C_8}{R}+\frac{C_9}{R^2}+C_{10}\right).
\]
Here \(C_8=C_8(n,K)\), \(C_9=C_9(n,\varepsilon)\), and
\(C_{10}=C_{10}(n,K,\varepsilon,C')\).
\end{theorem}
\begin{proof}
\sketch Apply Theorem~\ref{thm:chow-iii-li-yau-evolving-metrics}'s
minimum-principle argument with \(f(x,\tau)=d_{g(\tau)}(x,p)\).  Outside
\(B_{g(\tau)}(p,K^{-1/2})\), Theorem
\ref{thm:chow-iii-changing-distances} gives, in the barrier sense,
\[
 (\partial_\tau-\Delta)f\geq-\frac53(n-1)\sqrt K,
 \qquad |\nabla f|=1.
\]
For \(\phi=\psi(f/R)\), these imply the cutoff inequality with right side
\(C_8/R+C_9/R^2\).  Since \(R\geq K^{-1/2}\) and \(\phi=1\) on the inner
ball, it holds everywhere relevant to a negative minimum.  The same
quadratic minimum-principle calculation gives the displayed lower bound.
At cut-locus points use the upper distance barrier from the changing-distance
argument; because \(\psi\) decreases and \(P<0\), it is a valid lower barrier
for \(\tau\phi P\).
\end{proof}

\begin{remark}[Locality of the Ricci-flow hypotheses]
\label{rem:chow-iii-li-yau-locality-ricci-flow}
\source{chow-et-al-ricci-flow-part-iii:page-0340}
\dcref{ch25:25.10}
\group{Li--Yau Estimates}

Theorem~\ref{thm:chow-iii-li-yau-ricci-flow} requires its geometric bounds
only on the indicated moving ball, unlike Theorem
\ref{thm:chow-iii-li-yau-evolving-metrics}.
\end{remark}

\section{Integrated Harnack inequalities}

\begin{corollary}[Classical-type Harnack inequality for evolving metrics]
\label{cor:chow-iii-classical-harnack-evolving-metrics}
\source{chow-et-al-ricci-flow-part-iii:page-0340,chow-et-al-ricci-flow-part-iii:page-0341}
\uses{thm:chow-iii-li-yau-evolving-metrics}
\dcref{ch25:25.12}
\group{Integrated Harnack Inequalities}


Assume \(\widetilde C_0^{-1}\widetilde g\leq g(\tau)leq
\widetilde C_0\widetilde g\).  Under the hypotheses of Theorem
\ref{thm:chow-iii-li-yau-evolving-metrics}, for
\(x_1,x_2\in\mathcal M\) and \(0<\tau_1<\tau_2\leq T\),
\[
 \frac{u(x_2,\tau_2)}{u(x_1,\tau_1)}
 \geq e^{-C_{11}(\tau_2-\tau_1)}
 \left(\frac{\tau_2}{\tau_1}\right)^{-n/(2-3\varepsilon)}
 \exp\left(-\frac{\widetilde C_0d_{\widetilde g}(x_1,x_2)^2}
 {4(1-\varepsilon)(\tau_2-\tau_1)}\right),
\]
where \(C_{11}=C_7+\sup_{\mathcal M\times[\tau_1,\tau_2]}Q\).
\end{corollary}
\begin{proof}
\sketch Let \(\gamma\) join \(x_1\) to \(x_2\).  Along it, the global differential
estimate and
\[
 (1-\varepsilon)|\nabla\log u|^2
 +\langle\nabla\log u,\dot\gamma\rangle
 \geq-\frac{|\dot\gamma|^2}{4(1-\varepsilon)}
\]
give
\[
 \log\frac{u(x_2,\tau_2)}{u(x_1,\tau_1)}
 \geq-\frac1{4(1-\varepsilon)}\int_{\tau_1}^{\tau_2}|\dot\gamma|_{g(\tau)}^2d\tau
 -\frac n{2-3\varepsilon}\log\frac{\tau_2}{\tau_1}
 -C_{11}(\tau_2-\tau_1).
\]
Take \(\gamma\) to be a constant-speed minimizing \(\widetilde g\)-geodesic.
Metric comparison bounds its energy by
\(\widetilde C_0d_{\widetilde g}(x_1,x_2)^2/(\tau_2-\tau_1)\).
Exponentiation proves the claim.
\end{proof}

\begin{corollary}[Classical-type Harnack inequality for Ricci flow]
\label{cor:chow-iii-classical-harnack-ricci-flow}
\source{chow-et-al-ricci-flow-part-iii:page-0341,chow-et-al-ricci-flow-part-iii:page-0342}
\uses{thm:chow-iii-li-yau-ricci-flow}
\dcref{ch25:25.13}
\group{Integrated Harnack Inequalities}


Under the hypotheses of Theorem~\ref{thm:chow-iii-li-yau-ricci-flow}, assume
\(\widetilde C_0^{-1}\widetilde g\leq g(\tau)\leq
\widetilde C_0\widetilde g\) on \(B_{g(\tau)}(p,R)\).  If a minimizing
\(\widetilde g\)-geodesic from \(x_1\) to \(x_2\) remains in those balls,
then
\[
 \frac{u(x_2,\tau_2)}{u(x_1,\tau_1)}
 \geq e^{-C_{12}(\tau_2-\tau_1)}
 \left(\frac{\tau_2}{\tau_1}\right)^{-n/(2-3\varepsilon)}
 \exp\left(-\frac{\widetilde C_0d_{\widetilde g}(x_1,x_2)^2}
 {4(1-\varepsilon)(\tau_2-\tau_1)}\right),
\]
where
\(C_{12}=C_8/R+C_9/R^2+C_{10}+\sup Q\) on the parabolic region.
\end{corollary}
\begin{proof}
\sketch Integrate Theorem~\ref{thm:chow-iii-li-yau-ricci-flow} along the prescribed
geodesic and complete the square exactly as in Corollary
\ref{cor:chow-iii-classical-harnack-evolving-metrics}.  The hypothesis keeps
the path inside the region where the differential estimate holds, and metric
comparison gives the same energy bound.
\end{proof}

\begin{remark}[Mean value consequences and adjoint kernels]
\label{rem:chow-iii-harnack-mean-value-adjoint-kernels}
\source{chow-et-al-ricci-flow-part-iii:page-0342}
\dcref{ch25:25.15}
\group{Integrated Harnack Inequalities}

Integrating the classical Harnack inequality over a ball produces spatial
mean value upper and lower bounds.  In conjunction with heat-kernel symmetry,
these estimates give corresponding bounds for adjoint heat kernels under
local hypotheses.
\end{remark}

\section{Evolution and localization of the Harnack quantity}

Let \(L=\log u\).  Then
\[
 \partial_\tau L=\Delta L+|\nabla L|^2-Q,
 \qquad P=\partial_\tau L-(1-\varepsilon)|\nabla L|^2+Q
 =\Delta L+\varepsilon|\nabla L|^2.
\]

\begin{remark}[Model values of the Harnack quantity]
\label{rem:chow-iii-model-harnack-quantity}
\source{chow-et-al-ricci-flow-part-iii:page-0343}
\dcref{ch25:25.16}
\group{Harnack Quantity}

For the Euclidean heat kernel and \(\varepsilon=0\), one has
\(\tau P=-n/2\).  For the heat equation on a static manifold with
nonnegative Ricci curvature, the Li--Yau estimate is \(P\geq-n/(2\tau)\).
\end{remark}

\begin{lemma}[Evolution equation for the Harnack quantity]
\label{lem:chow-iii-harnack-quantity-evolution}
\source{chow-et-al-ricci-flow-part-iii:page-0343,chow-et-al-ricci-flow-part-iii:page-0344}
\dcref{ch25:25.17}
\group{Harnack Quantity}

For every \(\varepsilon\in\mathbb R\),
\[
\begin{aligned}
 \partial_\tau P={}&\Delta P+2\langle\nabla L,\nabla P\rangle
 +2(1-\varepsilon)|\Hess L|^2-2\mathcal R_{ij}\nabla_i\nabla_jL\\
 &+2\bigl((1-\varepsilon)R_{ij}-\varepsilon\mathcal R_{ij}\bigr)
   \nabla_iL\nabla_jL\\
 &+\bigl(-2g^{ij}\nabla_i\mathcal R_{jk}+\nabla_k\mathcal R
          -2\varepsilon\nabla_kQ\bigr)\nabla_kL-\Delta Q.
\end{aligned}
\]
\end{lemma}
\begin{proof}
\sketch The variation of the Laplacian is
\[
 \partial_\tau\Delta=-2\mathcal R_{ij}\nabla_i\nabla_j
 +\bigl(-2g^{ij}\nabla_i\mathcal R_{jk}+\nabla_k\mathcal R\bigr)\nabla_k.
\]
Differentiate \(P=\Delta L+\varepsilon|\nabla L|^2\), use
\(\partial_\tau L=P+(1-\varepsilon)|\nabla L|^2-Q\), and account for
\(\partial_\tau g^{-1}=-2\mathcal R^{\sharp}\).  Finally apply Bochner's
identity
\[
 \Delta|\nabla L|^2=2|\Hess L|^2
 +2\langle\nabla L,\nabla\Delta L\rangle
 +2\Ric(\nabla L,\nabla L).
\]
The terms differentiating \(|\nabla L|^2\) combine because
\(\nabla P=\nabla\Delta L+\varepsilon\nabla|\nabla L|^2\); collecting the
remaining terms gives the formula.
\end{proof}

\begin{corollary}[Evolution inequality for the Harnack quantity]
\label{cor:chow-iii-harnack-quantity-evolution-inequality}
\source{chow-et-al-ricci-flow-part-iii:page-0344,chow-et-al-ricci-flow-part-iii:page-0345}
\uses{lem:chow-iii-harnack-quantity-evolution}
\dcref{ch25:25.18}
\group{Harnack Quantity}


Suppose
\[
 \Ric\geq-K,\quad |\mathcal R_{ij}|\leq\Lambda,
\quad
 \left|-2g^{ij}\nabla_i\mathcal R_{jk}+\nabla_k\mathcal R\right|
 +\frac43|\nabla Q|\leq C',\quad \Delta Q\leq C'.
\]
If \(0<\varepsilon<2/3\), then
\[
 \partial_\tau P\geq\Delta P+2\langle\nabla L,\nabla P\rangle
 +\frac{2-3\varepsilon}{n}(\Delta L)^2
 -C_1|\nabla L|^2-C_2,
\]
where
\[
 C_1=2\bigl((1-\varepsilon)K+\varepsilon\Lambda+C'\bigr),
 \qquad C_2=\frac{\Lambda^2}{\varepsilon}+2C'.
\]
\end{corollary}
\begin{proof}
\sketch Insert the assumed bounds into Lemma
\ref{lem:chow-iii-harnack-quantity-evolution}.  Peter--Paul gives
\[
 2(1-\varepsilon)|\Hess L|^2-2\Lambda|\Hess L|
 \geq(2-3\varepsilon)|\Hess L|^2-\frac{\Lambda^2}{\varepsilon}
\]
and
\(-C'|\nabla L|\geq-C'|\nabla L|^2-C'\).  The curvature-gradient term is
bounded below by
\(-2((1-\varepsilon)K+\varepsilon\Lambda)|\nabla L|^2\).
Combining these estimates and using
\(|\Hess L|^2\geq(\Delta L)^2/n\) proves the inequality.
\end{proof}

\begin{example}[Exercise 25.14: integrate the Ricci-flow gradient estimate]
\label{ex:chow-iii-integrate-ricci-flow-gradient-estimate}
\dcref{ch25:25.14}
\group{Integrated Harnack Inequalities}
\source{chow-et-al-ricci-flow-part-iii:page-0342}
\uses{thm:chow-iii-li-yau-ricci-flow}
Derive Corollary~\ref{cor:chow-iii-classical-harnack-ricci-flow} from the
differential estimate.
\end{example}
\begin{proof}
Along a path \(\gamma\), differentiate \(\log u(\gamma(\tau),\tau)\), insert
the differential Harnack estimate, and apply
\(a|X|^2+\langle X,Y\rangle\geq-|Y|^2/(4a)\) with
\(a=1-\varepsilon\).  Integrating \(1/\tau\) gives the time-ratio factor;
taking the constant-speed minimizing \(\widetilde g\)-geodesic and applying
metric comparison gives the Gaussian factor.
\end{proof}

\begin{example}[Problem 25.11: local estimate for backward Ricci flow]
\label{ex:chow-iii-local-li-yau-backward-ricci-flow}
\dcref{ch25:25.11}
\group{Li--Yau Problems}
\source{chow-et-al-ricci-flow-part-iii:page-0340}
Find a Li--Yau differential Harnack estimate for positive solutions of
heat-type equations coupled to backward Ricci flow using only local bounds on
the curvature and potential.
\end{example}
